\documentclass[11pt,a4paper]{article}
\usepackage[utf8]{inputenc}
\usepackage[portuguese]{babel}
\usepackage[T1]{fontenc}
\usepackage[ddmmyyyy]{datetime}
\usepackage[margin=2cm]{geometry}
\usepackage{amsfonts}
\usepackage{amsmath}
\usepackage{amssymb}
\usepackage{amsthm}
\usepackage{biblatex}
\usepackage{csquotes}
\usepackage{graphicx}
\usepackage{mathtools}
\usepackage{thmtools}
\usepackage{titling}
\usepackage{xcolor}
\usepackage{enumitem}
\usepackage[unicode]{hyperref}

\pdfsuppressptexinfo=7

\graphicspath{{e-folio/e-folio-a/imagem}}
\addbibresource{bibliografia.bib}

% Ambientes de provas e demonstrações
\newtheorem{proposition}{Proposição}[section]
\renewcommand\qedsymbol{\textbf{q.e.d.}}

% Dados
\author{Carlos Augusto Gonçalves Collaço e Pinto Machado}
\newcommand{\numeroEstudante}{2200909}
\newcommand{\UC}{Elementos de Análise Infinitesimal I}
\newcommand{\codigoUC}{21030}
\newcommand{\docentes}{Maria João Oliveira}
\newcommand{\curso}{Licenciatura em Matemática e Aplicações}
\newcommand{\turma}{2}
\newcommand{\anoLectivo}{2023-2024}
\date{\today}

\pagenumbering{arabic}

\hypersetup{
	pdfsubject={\UC},
	pdftitle={E-fólio A - \numeroEstudante},
	pdfauthor={Carlos Pinto Machado},
	pdfproducer={},
	pdfcreator={}
}

\makeindex

\newcommand{\exercicio}{
	\addtocounter{section}{1}
	\setcounter{subsection}{0}
	\section*{\arabic{section}}
	\addcontentsline{toc}{section}{Exercício \arabic{section}}
}


\begin{document}

\includegraphics[width=\textwidth]{e-folio-a.jpg}

\paragraph{\textbf{UNIDADE CURRICULAR:}} \UC

\paragraph{\textbf{CÓDIGO:}} \codigoUC

\paragraph{\textbf{DOCENTE:}} \docentes

\paragraph{\textbf{NOME:}} \theauthor

\paragraph{\textbf{N.º DE ESTUDANTE:}} \numeroEstudante

\paragraph{\textbf{CURSO:}} \curso

\paragraph{\textbf{DATA DE ENTREGA:}} \thedate


\clearpage

\exercicio{}

\begin{equation}
	\begin{cases}
		v_1 = 1 \\
		v_{n + 1} = \frac{v^2_n}{1 + v_n}, \quad n \in \mathbb{N}
	\end{cases}
\end{equation}

\begin{enumerate}[label=\arabic{section}.\arabic*.]
	\item
	      \begin{proposition}
		      Para qualquer $n \in \mathbb{N}$, tem-se que $0 < v_n\leq 1$.
	      \end{proposition}
	      \begin{proof}
		      \hfill\\
		      \begin{enumerate}[label=\arabic*.]
			      \item Começamos pelo caso $n = 1$: temos $v_1 = 1 \in]0, 1]$
			            , que é uma afirmação verdadeira.
			      \item Fixamos agora um $n \in \mathbb{N}$ e assumimos que
			            \begin{align*}
				                 & 0 < v_n \leq 1 \quad\quad \text{(Hipótese de Indução)}
				            \intertext{Queremos mostrar que:}
				                 & 0 < v_{n + 1} \leq 1 \quad \text{(Tese de Indução)}
				            \intertext{Pela Hipótese de Indução, tem-se que}
				                 & 0 < v_n \leq 1 \implies 0 < v_n^2 \leq 1               \\
				                 & 0 < v_n \leq 1 \iff \; 1 < 1 + v_n
				            \intertext{Que nos permite, deduzir que:}
				                 & 0 < v_n^2 \leq 1 < 1 + v_n                             \\
				            \iff & 0 < \frac{v_n^2}{1 + v_n} \leq \frac{1}{1 + v_n} < 1
				            \intertext{Dado que
					            $v_{n + 1} = \frac{v_n^2}{1 + v_n}$, podemos
					            concluir que está demonstrado, por indução, que}
				                 & 0 < v_{n+1} \leq 1.
			            \end{align*}
		      \end{enumerate}
	      \end{proof}
	\item
	      \begin{proposition}
		      A sucessão $v_n$ é monótona.
	      \end{proposition}
	\item \hfill
	      \begin{enumerate}[label=\arabic{section}.\arabic{enumi}.\arabic*.]
		      \item
		            \begin{align*}
			            \sum_{n=1}^{\infty} v_n
		            \end{align*}
		      \item
		            \begin{align*}
			            \sum_{n=1}^{\infty} \frac{v_n}{1 + v_n}
		            \end{align*}
	      \end{enumerate}
\end{enumerate}

\exercicio{}

\begin{align*}
	\sum_{n=0}^{\infty} \frac{1}{1 + \alpha^{2n}}
\end{align*}

\begin{align*}
	(1 + \alpha^{2n})^{-1}
\end{align*}

\clearpage

\nocite{*}
\printbibliography[heading=bibintoc,title={Bibliografia}]

\end{document}
